\documentclass[11pt]{article}

    \usepackage[breakable]{tcolorbox}
    \usepackage{parskip} % Stop auto-indenting (to mimic markdown behaviour)
    

    % Basic figure setup, for now with no caption control since it's done
    % automatically by Pandoc (which extracts ![](path) syntax from Markdown).
    \usepackage{graphicx}
    % Keep aspect ratio if custom image width or height is specified
    \setkeys{Gin}{keepaspectratio}
    % Maintain compatibility with old templates. Remove in nbconvert 6.0
    \let\Oldincludegraphics\includegraphics
    % Ensure that by default, figures have no caption (until we provide a
    % proper Figure object with a Caption API and a way to capture that
    % in the conversion process - todo).
    \usepackage{caption}
    \DeclareCaptionFormat{nocaption}{}
    \captionsetup{format=nocaption,aboveskip=0pt,belowskip=0pt}

    \usepackage{float}
    \floatplacement{figure}{H} % forces figures to be placed at the correct location
    \usepackage{xcolor} % Allow colors to be defined
    \usepackage{enumerate} % Needed for markdown enumerations to work
    \usepackage{geometry} % Used to adjust the document margins
    \usepackage{amsmath} % Equations
    \usepackage{amssymb} % Equations
    \usepackage{textcomp} % defines textquotesingle
    % Hack from http://tex.stackexchange.com/a/47451/13684:
    \AtBeginDocument{%
        \def\PYZsq{\textquotesingle}% Upright quotes in Pygmentized code
    }
    \usepackage{upquote} % Upright quotes for verbatim code
    \usepackage{eurosym} % defines \euro

    \usepackage{iftex}
    \ifPDFTeX
        \usepackage[T1]{fontenc}
        \IfFileExists{alphabeta.sty}{
              \usepackage{alphabeta}
          }{
              \usepackage[mathletters]{ucs}
              \usepackage[utf8x]{inputenc}
          }
    \else
        \usepackage{fontspec}
        \usepackage{unicode-math}
    \fi

    \usepackage{fancyvrb} % verbatim replacement that allows latex
    \usepackage{grffile} % extends the file name processing of package graphics
                         % to support a larger range
    \makeatletter % fix for old versions of grffile with XeLaTeX
    \@ifpackagelater{grffile}{2019/11/01}
    {
      % Do nothing on new versions
    }
    {
      \def\Gread@@xetex#1{%
        \IfFileExists{"\Gin@base".bb}%
        {\Gread@eps{\Gin@base.bb}}%
        {\Gread@@xetex@aux#1}%
      }
    }
    \makeatother
    \usepackage[Export]{adjustbox} % Used to constrain images to a maximum size
    \adjustboxset{max size={0.9\linewidth}{0.9\paperheight}}

    % The hyperref package gives us a pdf with properly built
    % internal navigation ('pdf bookmarks' for the table of contents,
    % internal cross-reference links, web links for URLs, etc.)
    \usepackage{hyperref}
    % The default LaTeX title has an obnoxious amount of whitespace. By default,
    % titling removes some of it. It also provides customization options.
    \usepackage{titling}
    \usepackage{longtable} % longtable support required by pandoc >1.10
    \usepackage{booktabs}  % table support for pandoc > 1.12.2
    \usepackage{array}     % table support for pandoc >= 2.11.3
    \usepackage{calc}      % table minipage width calculation for pandoc >= 2.11.1
    \usepackage[inline]{enumitem} % IRkernel/repr support (it uses the enumerate* environment)
    \usepackage[normalem]{ulem} % ulem is needed to support strikethroughs (\sout)
                                % normalem makes italics be italics, not underlines
    \usepackage{soul}      % strikethrough (\st) support for pandoc >= 3.0.0
    \usepackage{mathrsfs}
    

    
    % Colors for the hyperref package
    \definecolor{urlcolor}{rgb}{0,.145,.698}
    \definecolor{linkcolor}{rgb}{.71,0.21,0.01}
    \definecolor{citecolor}{rgb}{.12,.54,.11}

    % ANSI colors
    \definecolor{ansi-black}{HTML}{3E424D}
    \definecolor{ansi-black-intense}{HTML}{282C36}
    \definecolor{ansi-red}{HTML}{E75C58}
    \definecolor{ansi-red-intense}{HTML}{B22B31}
    \definecolor{ansi-green}{HTML}{00A250}
    \definecolor{ansi-green-intense}{HTML}{007427}
    \definecolor{ansi-yellow}{HTML}{DDB62B}
    \definecolor{ansi-yellow-intense}{HTML}{B27D12}
    \definecolor{ansi-blue}{HTML}{208FFB}
    \definecolor{ansi-blue-intense}{HTML}{0065CA}
    \definecolor{ansi-magenta}{HTML}{D160C4}
    \definecolor{ansi-magenta-intense}{HTML}{A03196}
    \definecolor{ansi-cyan}{HTML}{60C6C8}
    \definecolor{ansi-cyan-intense}{HTML}{258F8F}
    \definecolor{ansi-white}{HTML}{C5C1B4}
    \definecolor{ansi-white-intense}{HTML}{A1A6B2}
    \definecolor{ansi-default-inverse-fg}{HTML}{FFFFFF}
    \definecolor{ansi-default-inverse-bg}{HTML}{000000}

    % common color for the border for error outputs.
    \definecolor{outerrorbackground}{HTML}{FFDFDF}

    % commands and environments needed by pandoc snippets
    % extracted from the output of `pandoc -s`
    \providecommand{\tightlist}{%
      \setlength{\itemsep}{0pt}\setlength{\parskip}{0pt}}
    \DefineVerbatimEnvironment{Highlighting}{Verbatim}{commandchars=\\\{\}}
    % Add ',fontsize=\small' for more characters per line
    \newenvironment{Shaded}{}{}
    \newcommand{\KeywordTok}[1]{\textcolor[rgb]{0.00,0.44,0.13}{\textbf{{#1}}}}
    \newcommand{\DataTypeTok}[1]{\textcolor[rgb]{0.56,0.13,0.00}{{#1}}}
    \newcommand{\DecValTok}[1]{\textcolor[rgb]{0.25,0.63,0.44}{{#1}}}
    \newcommand{\BaseNTok}[1]{\textcolor[rgb]{0.25,0.63,0.44}{{#1}}}
    \newcommand{\FloatTok}[1]{\textcolor[rgb]{0.25,0.63,0.44}{{#1}}}
    \newcommand{\CharTok}[1]{\textcolor[rgb]{0.25,0.44,0.63}{{#1}}}
    \newcommand{\StringTok}[1]{\textcolor[rgb]{0.25,0.44,0.63}{{#1}}}
    \newcommand{\CommentTok}[1]{\textcolor[rgb]{0.38,0.63,0.69}{\textit{{#1}}}}
    \newcommand{\OtherTok}[1]{\textcolor[rgb]{0.00,0.44,0.13}{{#1}}}
    \newcommand{\AlertTok}[1]{\textcolor[rgb]{1.00,0.00,0.00}{\textbf{{#1}}}}
    \newcommand{\FunctionTok}[1]{\textcolor[rgb]{0.02,0.16,0.49}{{#1}}}
    \newcommand{\RegionMarkerTok}[1]{{#1}}
    \newcommand{\ErrorTok}[1]{\textcolor[rgb]{1.00,0.00,0.00}{\textbf{{#1}}}}
    \newcommand{\NormalTok}[1]{{#1}}

    % Additional commands for more recent versions of Pandoc
    \newcommand{\ConstantTok}[1]{\textcolor[rgb]{0.53,0.00,0.00}{{#1}}}
    \newcommand{\SpecialCharTok}[1]{\textcolor[rgb]{0.25,0.44,0.63}{{#1}}}
    \newcommand{\VerbatimStringTok}[1]{\textcolor[rgb]{0.25,0.44,0.63}{{#1}}}
    \newcommand{\SpecialStringTok}[1]{\textcolor[rgb]{0.73,0.40,0.53}{{#1}}}
    \newcommand{\ImportTok}[1]{{#1}}
    \newcommand{\DocumentationTok}[1]{\textcolor[rgb]{0.73,0.13,0.13}{\textit{{#1}}}}
    \newcommand{\AnnotationTok}[1]{\textcolor[rgb]{0.38,0.63,0.69}{\textbf{\textit{{#1}}}}}
    \newcommand{\CommentVarTok}[1]{\textcolor[rgb]{0.38,0.63,0.69}{\textbf{\textit{{#1}}}}}
    \newcommand{\VariableTok}[1]{\textcolor[rgb]{0.10,0.09,0.49}{{#1}}}
    \newcommand{\ControlFlowTok}[1]{\textcolor[rgb]{0.00,0.44,0.13}{\textbf{{#1}}}}
    \newcommand{\OperatorTok}[1]{\textcolor[rgb]{0.40,0.40,0.40}{{#1}}}
    \newcommand{\BuiltInTok}[1]{{#1}}
    \newcommand{\ExtensionTok}[1]{{#1}}
    \newcommand{\PreprocessorTok}[1]{\textcolor[rgb]{0.74,0.48,0.00}{{#1}}}
    \newcommand{\AttributeTok}[1]{\textcolor[rgb]{0.49,0.56,0.16}{{#1}}}
    \newcommand{\InformationTok}[1]{\textcolor[rgb]{0.38,0.63,0.69}{\textbf{\textit{{#1}}}}}
    \newcommand{\WarningTok}[1]{\textcolor[rgb]{0.38,0.63,0.69}{\textbf{\textit{{#1}}}}}
    \makeatletter
    \newsavebox\pandoc@box
    \newcommand*\pandocbounded[1]{%
      \sbox\pandoc@box{#1}%
      % scaling factors for width and height
      \Gscale@div\@tempa\textheight{\dimexpr\ht\pandoc@box+\dp\pandoc@box\relax}%
      \Gscale@div\@tempb\linewidth{\wd\pandoc@box}%
      % select the smaller of both
      \ifdim\@tempb\p@<\@tempa\p@
        \let\@tempa\@tempb
      \fi
      % scaling accordingly (\@tempa < 1)
      \ifdim\@tempa\p@<\p@
        \scalebox{\@tempa}{\usebox\pandoc@box}%
      % scaling not needed, use as it is
      \else
        \usebox{\pandoc@box}%
      \fi
    }
    \makeatother

    % Define a nice break command that doesn't care if a line doesn't already
    % exist.
    \def\br{\hspace*{\fill} \\* }
    % Math Jax compatibility definitions
    \def\gt{>}
    \def\lt{<}
    \let\Oldtex\TeX
    \let\Oldlatex\LaTeX
    \renewcommand{\TeX}{\textrm{\Oldtex}}
    \renewcommand{\LaTeX}{\textrm{\Oldlatex}}
    % Document parameters
    % Document title
    \title{open\_data\_collaborator\_code\_notebook}
    
    
    
    
    
    
    
% Pygments definitions
\makeatletter
\def\PY@reset{\let\PY@it=\relax \let\PY@bf=\relax%
    \let\PY@ul=\relax \let\PY@tc=\relax%
    \let\PY@bc=\relax \let\PY@ff=\relax}
\def\PY@tok#1{\csname PY@tok@#1\endcsname}
\def\PY@toks#1+{\ifx\relax#1\empty\else%
    \PY@tok{#1}\expandafter\PY@toks\fi}
\def\PY@do#1{\PY@bc{\PY@tc{\PY@ul{%
    \PY@it{\PY@bf{\PY@ff{#1}}}}}}}
\def\PY#1#2{\PY@reset\PY@toks#1+\relax+\PY@do{#2}}

\@namedef{PY@tok@w}{\def\PY@tc##1{\textcolor[rgb]{0.73,0.73,0.73}{##1}}}
\@namedef{PY@tok@c}{\let\PY@it=\textit\def\PY@tc##1{\textcolor[rgb]{0.24,0.48,0.48}{##1}}}
\@namedef{PY@tok@cp}{\def\PY@tc##1{\textcolor[rgb]{0.61,0.40,0.00}{##1}}}
\@namedef{PY@tok@k}{\let\PY@bf=\textbf\def\PY@tc##1{\textcolor[rgb]{0.00,0.50,0.00}{##1}}}
\@namedef{PY@tok@kp}{\def\PY@tc##1{\textcolor[rgb]{0.00,0.50,0.00}{##1}}}
\@namedef{PY@tok@kt}{\def\PY@tc##1{\textcolor[rgb]{0.69,0.00,0.25}{##1}}}
\@namedef{PY@tok@o}{\def\PY@tc##1{\textcolor[rgb]{0.40,0.40,0.40}{##1}}}
\@namedef{PY@tok@ow}{\let\PY@bf=\textbf\def\PY@tc##1{\textcolor[rgb]{0.67,0.13,1.00}{##1}}}
\@namedef{PY@tok@nb}{\def\PY@tc##1{\textcolor[rgb]{0.00,0.50,0.00}{##1}}}
\@namedef{PY@tok@nf}{\def\PY@tc##1{\textcolor[rgb]{0.00,0.00,1.00}{##1}}}
\@namedef{PY@tok@nc}{\let\PY@bf=\textbf\def\PY@tc##1{\textcolor[rgb]{0.00,0.00,1.00}{##1}}}
\@namedef{PY@tok@nn}{\let\PY@bf=\textbf\def\PY@tc##1{\textcolor[rgb]{0.00,0.00,1.00}{##1}}}
\@namedef{PY@tok@ne}{\let\PY@bf=\textbf\def\PY@tc##1{\textcolor[rgb]{0.80,0.25,0.22}{##1}}}
\@namedef{PY@tok@nv}{\def\PY@tc##1{\textcolor[rgb]{0.10,0.09,0.49}{##1}}}
\@namedef{PY@tok@no}{\def\PY@tc##1{\textcolor[rgb]{0.53,0.00,0.00}{##1}}}
\@namedef{PY@tok@nl}{\def\PY@tc##1{\textcolor[rgb]{0.46,0.46,0.00}{##1}}}
\@namedef{PY@tok@ni}{\let\PY@bf=\textbf\def\PY@tc##1{\textcolor[rgb]{0.44,0.44,0.44}{##1}}}
\@namedef{PY@tok@na}{\def\PY@tc##1{\textcolor[rgb]{0.41,0.47,0.13}{##1}}}
\@namedef{PY@tok@nt}{\let\PY@bf=\textbf\def\PY@tc##1{\textcolor[rgb]{0.00,0.50,0.00}{##1}}}
\@namedef{PY@tok@nd}{\def\PY@tc##1{\textcolor[rgb]{0.67,0.13,1.00}{##1}}}
\@namedef{PY@tok@s}{\def\PY@tc##1{\textcolor[rgb]{0.73,0.13,0.13}{##1}}}
\@namedef{PY@tok@sd}{\let\PY@it=\textit\def\PY@tc##1{\textcolor[rgb]{0.73,0.13,0.13}{##1}}}
\@namedef{PY@tok@si}{\let\PY@bf=\textbf\def\PY@tc##1{\textcolor[rgb]{0.64,0.35,0.47}{##1}}}
\@namedef{PY@tok@se}{\let\PY@bf=\textbf\def\PY@tc##1{\textcolor[rgb]{0.67,0.36,0.12}{##1}}}
\@namedef{PY@tok@sr}{\def\PY@tc##1{\textcolor[rgb]{0.64,0.35,0.47}{##1}}}
\@namedef{PY@tok@ss}{\def\PY@tc##1{\textcolor[rgb]{0.10,0.09,0.49}{##1}}}
\@namedef{PY@tok@sx}{\def\PY@tc##1{\textcolor[rgb]{0.00,0.50,0.00}{##1}}}
\@namedef{PY@tok@m}{\def\PY@tc##1{\textcolor[rgb]{0.40,0.40,0.40}{##1}}}
\@namedef{PY@tok@gh}{\let\PY@bf=\textbf\def\PY@tc##1{\textcolor[rgb]{0.00,0.00,0.50}{##1}}}
\@namedef{PY@tok@gu}{\let\PY@bf=\textbf\def\PY@tc##1{\textcolor[rgb]{0.50,0.00,0.50}{##1}}}
\@namedef{PY@tok@gd}{\def\PY@tc##1{\textcolor[rgb]{0.63,0.00,0.00}{##1}}}
\@namedef{PY@tok@gi}{\def\PY@tc##1{\textcolor[rgb]{0.00,0.52,0.00}{##1}}}
\@namedef{PY@tok@gr}{\def\PY@tc##1{\textcolor[rgb]{0.89,0.00,0.00}{##1}}}
\@namedef{PY@tok@ge}{\let\PY@it=\textit}
\@namedef{PY@tok@gs}{\let\PY@bf=\textbf}
\@namedef{PY@tok@ges}{\let\PY@bf=\textbf\let\PY@it=\textit}
\@namedef{PY@tok@gp}{\let\PY@bf=\textbf\def\PY@tc##1{\textcolor[rgb]{0.00,0.00,0.50}{##1}}}
\@namedef{PY@tok@go}{\def\PY@tc##1{\textcolor[rgb]{0.44,0.44,0.44}{##1}}}
\@namedef{PY@tok@gt}{\def\PY@tc##1{\textcolor[rgb]{0.00,0.27,0.87}{##1}}}
\@namedef{PY@tok@err}{\def\PY@bc##1{{\setlength{\fboxsep}{\string -\fboxrule}\fcolorbox[rgb]{1.00,0.00,0.00}{1,1,1}{\strut ##1}}}}
\@namedef{PY@tok@kc}{\let\PY@bf=\textbf\def\PY@tc##1{\textcolor[rgb]{0.00,0.50,0.00}{##1}}}
\@namedef{PY@tok@kd}{\let\PY@bf=\textbf\def\PY@tc##1{\textcolor[rgb]{0.00,0.50,0.00}{##1}}}
\@namedef{PY@tok@kn}{\let\PY@bf=\textbf\def\PY@tc##1{\textcolor[rgb]{0.00,0.50,0.00}{##1}}}
\@namedef{PY@tok@kr}{\let\PY@bf=\textbf\def\PY@tc##1{\textcolor[rgb]{0.00,0.50,0.00}{##1}}}
\@namedef{PY@tok@bp}{\def\PY@tc##1{\textcolor[rgb]{0.00,0.50,0.00}{##1}}}
\@namedef{PY@tok@fm}{\def\PY@tc##1{\textcolor[rgb]{0.00,0.00,1.00}{##1}}}
\@namedef{PY@tok@vc}{\def\PY@tc##1{\textcolor[rgb]{0.10,0.09,0.49}{##1}}}
\@namedef{PY@tok@vg}{\def\PY@tc##1{\textcolor[rgb]{0.10,0.09,0.49}{##1}}}
\@namedef{PY@tok@vi}{\def\PY@tc##1{\textcolor[rgb]{0.10,0.09,0.49}{##1}}}
\@namedef{PY@tok@vm}{\def\PY@tc##1{\textcolor[rgb]{0.10,0.09,0.49}{##1}}}
\@namedef{PY@tok@sa}{\def\PY@tc##1{\textcolor[rgb]{0.73,0.13,0.13}{##1}}}
\@namedef{PY@tok@sb}{\def\PY@tc##1{\textcolor[rgb]{0.73,0.13,0.13}{##1}}}
\@namedef{PY@tok@sc}{\def\PY@tc##1{\textcolor[rgb]{0.73,0.13,0.13}{##1}}}
\@namedef{PY@tok@dl}{\def\PY@tc##1{\textcolor[rgb]{0.73,0.13,0.13}{##1}}}
\@namedef{PY@tok@s2}{\def\PY@tc##1{\textcolor[rgb]{0.73,0.13,0.13}{##1}}}
\@namedef{PY@tok@sh}{\def\PY@tc##1{\textcolor[rgb]{0.73,0.13,0.13}{##1}}}
\@namedef{PY@tok@s1}{\def\PY@tc##1{\textcolor[rgb]{0.73,0.13,0.13}{##1}}}
\@namedef{PY@tok@mb}{\def\PY@tc##1{\textcolor[rgb]{0.40,0.40,0.40}{##1}}}
\@namedef{PY@tok@mf}{\def\PY@tc##1{\textcolor[rgb]{0.40,0.40,0.40}{##1}}}
\@namedef{PY@tok@mh}{\def\PY@tc##1{\textcolor[rgb]{0.40,0.40,0.40}{##1}}}
\@namedef{PY@tok@mi}{\def\PY@tc##1{\textcolor[rgb]{0.40,0.40,0.40}{##1}}}
\@namedef{PY@tok@il}{\def\PY@tc##1{\textcolor[rgb]{0.40,0.40,0.40}{##1}}}
\@namedef{PY@tok@mo}{\def\PY@tc##1{\textcolor[rgb]{0.40,0.40,0.40}{##1}}}
\@namedef{PY@tok@ch}{\let\PY@it=\textit\def\PY@tc##1{\textcolor[rgb]{0.24,0.48,0.48}{##1}}}
\@namedef{PY@tok@cm}{\let\PY@it=\textit\def\PY@tc##1{\textcolor[rgb]{0.24,0.48,0.48}{##1}}}
\@namedef{PY@tok@cpf}{\let\PY@it=\textit\def\PY@tc##1{\textcolor[rgb]{0.24,0.48,0.48}{##1}}}
\@namedef{PY@tok@c1}{\let\PY@it=\textit\def\PY@tc##1{\textcolor[rgb]{0.24,0.48,0.48}{##1}}}
\@namedef{PY@tok@cs}{\let\PY@it=\textit\def\PY@tc##1{\textcolor[rgb]{0.24,0.48,0.48}{##1}}}

\def\PYZbs{\char`\\}
\def\PYZus{\char`\_}
\def\PYZob{\char`\{}
\def\PYZcb{\char`\}}
\def\PYZca{\char`\^}
\def\PYZam{\char`\&}
\def\PYZlt{\char`\<}
\def\PYZgt{\char`\>}
\def\PYZsh{\char`\#}
\def\PYZpc{\char`\%}
\def\PYZdl{\char`\$}
\def\PYZhy{\char`\-}
\def\PYZsq{\char`\'}
\def\PYZdq{\char`\"}
\def\PYZti{\char`\~}
% for compatibility with earlier versions
\def\PYZat{@}
\def\PYZlb{[}
\def\PYZrb{]}
\makeatother


    % For linebreaks inside Verbatim environment from package fancyvrb.
    \makeatletter
        \newbox\Wrappedcontinuationbox
        \newbox\Wrappedvisiblespacebox
        \newcommand*\Wrappedvisiblespace {\textcolor{red}{\textvisiblespace}}
        \newcommand*\Wrappedcontinuationsymbol {\textcolor{red}{\llap{\tiny$\m@th\hookrightarrow$}}}
        \newcommand*\Wrappedcontinuationindent {3ex }
        \newcommand*\Wrappedafterbreak {\kern\Wrappedcontinuationindent\copy\Wrappedcontinuationbox}
        % Take advantage of the already applied Pygments mark-up to insert
        % potential linebreaks for TeX processing.
        %        {, <, #, %, $, ' and ": go to next line.
        %        _, }, ^, &, >, - and ~: stay at end of broken line.
        % Use of \textquotesingle for straight quote.
        \newcommand*\Wrappedbreaksatspecials {%
            \def\PYGZus{\discretionary{\char`\_}{\Wrappedafterbreak}{\char`\_}}%
            \def\PYGZob{\discretionary{}{\Wrappedafterbreak\char`\{}{\char`\{}}%
            \def\PYGZcb{\discretionary{\char`\}}{\Wrappedafterbreak}{\char`\}}}%
            \def\PYGZca{\discretionary{\char`\^}{\Wrappedafterbreak}{\char`\^}}%
            \def\PYGZam{\discretionary{\char`\&}{\Wrappedafterbreak}{\char`\&}}%
            \def\PYGZlt{\discretionary{}{\Wrappedafterbreak\char`\<}{\char`\<}}%
            \def\PYGZgt{\discretionary{\char`\>}{\Wrappedafterbreak}{\char`\>}}%
            \def\PYGZsh{\discretionary{}{\Wrappedafterbreak\char`\#}{\char`\#}}%
            \def\PYGZpc{\discretionary{}{\Wrappedafterbreak\char`\%}{\char`\%}}%
            \def\PYGZdl{\discretionary{}{\Wrappedafterbreak\char`\$}{\char`\$}}%
            \def\PYGZhy{\discretionary{\char`\-}{\Wrappedafterbreak}{\char`\-}}%
            \def\PYGZsq{\discretionary{}{\Wrappedafterbreak\textquotesingle}{\textquotesingle}}%
            \def\PYGZdq{\discretionary{}{\Wrappedafterbreak\char`\"}{\char`\"}}%
            \def\PYGZti{\discretionary{\char`\~}{\Wrappedafterbreak}{\char`\~}}%
        }
        % Some characters . , ; ? ! / are not pygmentized.
        % This macro makes them "active" and they will insert potential linebreaks
        \newcommand*\Wrappedbreaksatpunct {%
            \lccode`\~`\.\lowercase{\def~}{\discretionary{\hbox{\char`\.}}{\Wrappedafterbreak}{\hbox{\char`\.}}}%
            \lccode`\~`\,\lowercase{\def~}{\discretionary{\hbox{\char`\,}}{\Wrappedafterbreak}{\hbox{\char`\,}}}%
            \lccode`\~`\;\lowercase{\def~}{\discretionary{\hbox{\char`\;}}{\Wrappedafterbreak}{\hbox{\char`\;}}}%
            \lccode`\~`\:\lowercase{\def~}{\discretionary{\hbox{\char`\:}}{\Wrappedafterbreak}{\hbox{\char`\:}}}%
            \lccode`\~`\?\lowercase{\def~}{\discretionary{\hbox{\char`\?}}{\Wrappedafterbreak}{\hbox{\char`\?}}}%
            \lccode`\~`\!\lowercase{\def~}{\discretionary{\hbox{\char`\!}}{\Wrappedafterbreak}{\hbox{\char`\!}}}%
            \lccode`\~`\/\lowercase{\def~}{\discretionary{\hbox{\char`\/}}{\Wrappedafterbreak}{\hbox{\char`\/}}}%
            \catcode`\.\active
            \catcode`\,\active
            \catcode`\;\active
            \catcode`\:\active
            \catcode`\?\active
            \catcode`\!\active
            \catcode`\/\active
            \lccode`\~`\~
        }
    \makeatother

    \let\OriginalVerbatim=\Verbatim
    \makeatletter
    \renewcommand{\Verbatim}[1][1]{%
        %\parskip\z@skip
        \sbox\Wrappedcontinuationbox {\Wrappedcontinuationsymbol}%
        \sbox\Wrappedvisiblespacebox {\FV@SetupFont\Wrappedvisiblespace}%
        \def\FancyVerbFormatLine ##1{\hsize\linewidth
            \vtop{\raggedright\hyphenpenalty\z@\exhyphenpenalty\z@
                \doublehyphendemerits\z@\finalhyphendemerits\z@
                \strut ##1\strut}%
        }%
        % If the linebreak is at a space, the latter will be displayed as visible
        % space at end of first line, and a continuation symbol starts next line.
        % Stretch/shrink are however usually zero for typewriter font.
        \def\FV@Space {%
            \nobreak\hskip\z@ plus\fontdimen3\font minus\fontdimen4\font
            \discretionary{\copy\Wrappedvisiblespacebox}{\Wrappedafterbreak}
            {\kern\fontdimen2\font}%
        }%

        % Allow breaks at special characters using \PYG... macros.
        \Wrappedbreaksatspecials
        % Breaks at punctuation characters . , ; ? ! and / need catcode=\active
        \OriginalVerbatim[#1,codes*=\Wrappedbreaksatpunct]%
    }
    \makeatother

    % Exact colors from NB
    \definecolor{incolor}{HTML}{303F9F}
    \definecolor{outcolor}{HTML}{D84315}
    \definecolor{cellborder}{HTML}{CFCFCF}
    \definecolor{cellbackground}{HTML}{F7F7F7}

    % prompt
    \makeatletter
    \newcommand{\boxspacing}{\kern\kvtcb@left@rule\kern\kvtcb@boxsep}
    \makeatother
    \newcommand{\prompt}[4]{
        {\ttfamily\llap{{\color{#2}[#3]:\hspace{3pt}#4}}\vspace{-\baselineskip}}
    }
    

    
    % Prevent overflowing lines due to hard-to-break entities
    \sloppy
    % Setup hyperref package
    \hypersetup{
      breaklinks=true,  % so long urls are correctly broken across lines
      colorlinks=true,
      urlcolor=urlcolor,
      linkcolor=linkcolor,
      citecolor=citecolor,
      }
    % Slightly bigger margins than the latex defaults
    
    \geometry{verbose,tmargin=1in,bmargin=1in,lmargin=1in,rmargin=1in}
    
    

\begin{document}
    
    \maketitle
    
    

    
    \section{Open Data Notebook for Collaborators}\label{open-data-notebook-for-collaborators}

This notebook is the workflow for the Delft3D-FM UGRID model outputs.

This notebook:

\begin{itemize}
\tightlist
\item
  open the shared NetCDF file
\item
  preview key variables
\item
  make simple figures
\item
  optionally create an animation from the exported data
\end{itemize}

    \subsection{Import packages}\label{import-packages}

This section loads the Python packages used in the notebook.

These imports help with:

\begin{itemize}
\tightlist
\item
  reading NetCDF files with \texttt{xarray}
\item
  plotting with \texttt{matplotlib}
\item
  using ocean-focused colormaps from \texttt{cmocean}
\item
  working with unstructured grid model output through \texttt{xugrid}
\item
  optional animation support for saving movies
\end{itemize}

If \texttt{xugrid} is not installed, the import cell will stop and show
a clear error message.

    \begin{tcolorbox}[breakable, size=fbox, boxrule=1pt, pad at break*=1mm,colback=cellbackground, colframe=cellborder]
\prompt{In}{incolor}{6}{\boxspacing}
\begin{Verbatim}[commandchars=\\\{\}]
\PY{c+c1}{\PYZsh{} Optional one\PYZhy{}time install for video export only.}
\PY{c+c1}{\PYZsh{} Keep this False for normal use.}
\PY{n}{INSTALL\PYZus{}VIDEO\PYZus{}DEPS} \PY{o}{=} \PY{k+kc}{False}
\PY{k}{if} \PY{n}{INSTALL\PYZus{}VIDEO\PYZus{}DEPS}\PY{p}{:}
    \PY{k+kn}{import}\PY{+w}{ }\PY{n+nn}{subprocess}
    \PY{k+kn}{import}\PY{+w}{ }\PY{n+nn}{sys}

    \PY{c+c1}{\PYZsh{} Needed so Matplotlib can find ffmpeg for MP4 writing.}
    \PY{n}{subprocess}\PY{o}{.}\PY{n}{check\PYZus{}call}\PY{p}{(}\PY{p}{[}\PY{n}{sys}\PY{o}{.}\PY{n}{executable}\PY{p}{,} \PY{l+s+s2}{\PYZdq{}}\PY{l+s+s2}{\PYZhy{}m}\PY{l+s+s2}{\PYZdq{}}\PY{p}{,} \PY{l+s+s2}{\PYZdq{}}\PY{l+s+s2}{pip}\PY{l+s+s2}{\PYZdq{}}\PY{p}{,} \PY{l+s+s2}{\PYZdq{}}\PY{l+s+s2}{install}\PY{l+s+s2}{\PYZdq{}}\PY{p}{,} \PY{l+s+s2}{\PYZdq{}}\PY{l+s+s2}{\PYZhy{}U}\PY{l+s+s2}{\PYZdq{}}\PY{p}{,} \PY{l+s+s2}{\PYZdq{}}\PY{l+s+s2}{imageio\PYZhy{}ffmpeg}\PY{l+s+s2}{\PYZdq{}}\PY{p}{]}\PY{p}{)}

\PY{k+kn}{from}\PY{+w}{ }\PY{n+nn}{matplotlib}\PY{+w}{ }\PY{k+kn}{import} \PY{n}{pyplot} \PY{k}{as} \PY{n}{plt}
\PY{k+kn}{import}\PY{+w}{ }\PY{n+nn}{xarray}\PY{+w}{ }\PY{k}{as}\PY{+w}{ }\PY{n+nn}{xr}
\PY{k+kn}{import}\PY{+w}{ }\PY{n+nn}{cmocean}
\PY{k+kn}{from}\PY{+w}{ }\PY{n+nn}{pathlib}\PY{+w}{ }\PY{k+kn}{import} \PY{n}{Path}
\PY{k+kn}{import}\PY{+w}{ }\PY{n+nn}{xugrid}\PY{+w}{ }\PY{k}{as}\PY{+w}{ }\PY{n+nn}{xu}
\end{Verbatim}
\end{tcolorbox}

    \subsection{Set file path}\label{set-file-path}

Set the path to the NetCDF file you want to review.

    \begin{tcolorbox}[breakable, size=fbox, boxrule=1pt, pad at break*=1mm,colback=cellbackground, colframe=cellborder]
\prompt{In}{incolor}{7}{\boxspacing}
\begin{Verbatim}[commandchars=\\\{\}]
\PY{c+c1}{\PYZsh{} Paste your full NetCDF path here:}
\PY{n}{model\PYZus{}nc} \PY{o}{=} \PY{n}{Path}\PY{p}{(}\PY{l+s+sa}{r}\PY{l+s+s2}{\PYZdq{}}\PY{l+s+s2}{c:}\PY{l+s+s2}{\PYZbs{}}\PY{l+s+s2}{Users}\PY{l+s+s2}{\PYZbs{}}\PY{l+s+s2}{hydro}\PY{l+s+s2}{\PYZbs{}}\PY{l+s+s2}{Lake\PYZus{}Pontchartrain\PYZus{}Model}\PY{l+s+s2}{\PYZbs{}}\PY{l+s+s2}{outputs}\PY{l+s+s2}{\PYZbs{}}\PY{l+s+s2}{shared}\PY{l+s+s2}{\PYZbs{}}\PY{l+s+s2}{lp\PYZus{}model\PYZus{}spatial\PYZus{}window\PYZus{}20130601\PYZus{}20131001\PYZus{}ugrid.nc}\PY{l+s+s2}{\PYZdq{}}\PY{p}{)}
\end{Verbatim}
\end{tcolorbox}

    \subsection{Rename Variables}\label{rename-variables}

This section renames the main model variables to simpler labels used in
the plots below.

It also reads the \texttt{time} coordinate and converts it to a readable
string format for figure titles and summaries. After this cell runs, you
should see the dataset name, the time range in the file, and the number
of timesteps available.

    \begin{tcolorbox}[breakable, size=fbox, boxrule=1pt, pad at break*=1mm,colback=cellbackground, colframe=cellborder]
\prompt{In}{incolor}{8}{\boxspacing}
\begin{Verbatim}[commandchars=\\\{\}]
\PY{n}{var\PYZus{}name\PYZus{}map} \PY{o}{=} \PY{p}{\PYZob{}}
    \PY{l+s+s2}{\PYZdq{}}\PY{l+s+s2}{mesh2d\PYZus{}ucmag}\PY{l+s+s2}{\PYZdq{}}\PY{p}{:} \PY{l+s+s2}{\PYZdq{}}\PY{l+s+s2}{velocity\PYZus{}magnitude}\PY{l+s+s2}{\PYZdq{}}\PY{p}{,}
    \PY{l+s+s2}{\PYZdq{}}\PY{l+s+s2}{mesh2d\PYZus{}s1}\PY{l+s+s2}{\PYZdq{}}\PY{p}{:} \PY{l+s+s2}{\PYZdq{}}\PY{l+s+s2}{water\PYZus{}level}\PY{l+s+s2}{\PYZdq{}}\PY{p}{,}
    \PY{l+s+s2}{\PYZdq{}}\PY{l+s+s2}{mesh2d\PYZus{}tem1}\PY{l+s+s2}{\PYZdq{}}\PY{p}{:} \PY{l+s+s2}{\PYZdq{}}\PY{l+s+s2}{temperature}\PY{l+s+s2}{\PYZdq{}}\PY{p}{,}
    \PY{l+s+s2}{\PYZdq{}}\PY{l+s+s2}{mesh2d\PYZus{}sa1}\PY{l+s+s2}{\PYZdq{}}\PY{p}{:} \PY{l+s+s2}{\PYZdq{}}\PY{l+s+s2}{salinity}\PY{l+s+s2}{\PYZdq{}}\PY{p}{,}
\PY{p}{\PYZcb{}}

\PY{c+c1}{\PYZsh{} Build readable timestamps for plot titles and summaries.}
\PY{k}{with} \PY{n}{xr}\PY{o}{.}\PY{n}{open\PYZus{}dataset}\PY{p}{(}\PY{n}{model\PYZus{}nc}\PY{p}{)} \PY{k}{as} \PY{n}{ds\PYZus{}selected}\PY{p}{:}
    \PY{n}{time\PYZus{}str\PYZus{}all} \PY{o}{=} \PY{n}{ds\PYZus{}selected}\PY{o}{.}\PY{n}{time}\PY{o}{.}\PY{n}{dt}\PY{o}{.}\PY{n}{strftime}\PY{p}{(}\PY{l+s+s2}{\PYZdq{}}\PY{l+s+s2}{\PYZpc{}}\PY{l+s+s2}{m\PYZhy{}}\PY{l+s+si}{\PYZpc{}d}\PY{l+s+s2}{\PYZhy{}}\PY{l+s+s2}{\PYZpc{}}\PY{l+s+s2}{Y }\PY{l+s+s2}{\PYZpc{}}\PY{l+s+s2}{H:}\PY{l+s+s2}{\PYZpc{}}\PY{l+s+s2}{M:}\PY{l+s+s2}{\PYZpc{}}\PY{l+s+s2}{S}\PY{l+s+s2}{\PYZdq{}}\PY{p}{)}

    \PY{n}{t\PYZus{}start} \PY{o}{=} \PY{n}{time\PYZus{}str\PYZus{}all}\PY{o}{.}\PY{n}{isel}\PY{p}{(}\PY{n}{time}\PY{o}{=}\PY{l+m+mi}{0}\PY{p}{)}\PY{o}{.}\PY{n}{item}\PY{p}{(}\PY{p}{)}
    \PY{n}{t\PYZus{}end} \PY{o}{=} \PY{n}{time\PYZus{}str\PYZus{}all}\PY{o}{.}\PY{n}{isel}\PY{p}{(}\PY{n}{time}\PY{o}{=}\PY{o}{\PYZhy{}}\PY{l+m+mi}{1}\PY{p}{)}\PY{o}{.}\PY{n}{item}\PY{p}{(}\PY{p}{)}
    \PY{n}{n\PYZus{}timesteps} \PY{o}{=} \PY{n}{ds\PYZus{}selected}\PY{o}{.}\PY{n}{sizes}\PY{p}{[}\PY{l+s+s2}{\PYZdq{}}\PY{l+s+s2}{time}\PY{l+s+s2}{\PYZdq{}}\PY{p}{]}

\PY{n+nb}{print}\PY{p}{(}\PY{l+s+s2}{\PYZdq{}}\PY{l+s+se}{\PYZbs{}n}\PY{l+s+s2}{Dataset:}\PY{l+s+s2}{\PYZdq{}}\PY{p}{,} \PY{n}{model\PYZus{}nc}\PY{o}{.}\PY{n}{name}\PY{p}{)}
\PY{n+nb}{print}\PY{p}{(}\PY{l+s+s2}{\PYZdq{}}\PY{l+s+s2}{Time range:}\PY{l+s+s2}{\PYZdq{}}\PY{p}{,} \PY{n}{t\PYZus{}start}\PY{p}{,} \PY{l+s+s2}{\PYZdq{}}\PY{l+s+s2}{to}\PY{l+s+s2}{\PYZdq{}}\PY{p}{,} \PY{n}{t\PYZus{}end}\PY{p}{)}
\PY{n+nb}{print}\PY{p}{(}\PY{l+s+s2}{\PYZdq{}}\PY{l+s+s2}{Timesteps:}\PY{l+s+s2}{\PYZdq{}}\PY{p}{,} \PY{n}{n\PYZus{}timesteps}\PY{p}{)}
\end{Verbatim}
\end{tcolorbox}

    \begin{Verbatim}[commandchars=\\\{\}]

Dataset: lp\_model\_spatial\_window\_20130601\_20131001\_ugrid.nc
Time range: 06-01-2013 00:00:00 to 10-01-2013 00:00:00
Timesteps: 2929
    \end{Verbatim}

    \subsection{Create subplot preview figures}\label{create-subplot-preview-figures}

This section makes a quick four-panel map for one timestep.

The goal is to check of the main variables in the
file:

\begin{itemize}
\tightlist
\item
  velocity magnitude
\item
  water level
\item
  temperature
\item
  salinity
\end{itemize}

Each panel uses a fixed colormap and value range so the figures are
easier to compare across runs or exports.

    \begin{tcolorbox}[breakable, size=fbox, boxrule=1pt, pad at break*=1mm,colback=cellbackground, colframe=cellborder]
\prompt{In}{incolor}{9}{\boxspacing}
\begin{Verbatim}[commandchars=\\\{\}]
\PY{c+c1}{\PYZsh{} Set per\PYZhy{}variable colormap and colorbar limits (vmin/vmax).}
\PY{n}{plot\PYZus{}specs} \PY{o}{=} \PY{p}{[}
    \PY{p}{(}\PY{l+s+s2}{\PYZdq{}}\PY{l+s+s2}{velocity\PYZus{}magnitude}\PY{l+s+s2}{\PYZdq{}}\PY{p}{,} \PY{l+s+s2}{\PYZdq{}}\PY{l+s+s2}{Velocity Magnitude \PYZdl{}(m s\PYZca{}}\PY{l+s+s2}{\PYZob{}}\PY{l+s+s2}{\PYZhy{}1\PYZcb{})\PYZdl{}}\PY{l+s+s2}{\PYZdq{}}\PY{p}{,} \PY{n}{cmocean}\PY{o}{.}\PY{n}{cm}\PY{o}{.}\PY{n}{amp}\PY{p}{,} \PY{l+m+mf}{0.0}\PY{p}{,} \PY{l+m+mf}{1.2}\PY{p}{)}\PY{p}{,}
    \PY{p}{(}\PY{l+s+s2}{\PYZdq{}}\PY{l+s+s2}{water\PYZus{}level}\PY{l+s+s2}{\PYZdq{}}\PY{p}{,} \PY{l+s+s2}{\PYZdq{}}\PY{l+s+s2}{Water Level (m)}\PY{l+s+s2}{\PYZdq{}}\PY{p}{,} \PY{n}{cmocean}\PY{o}{.}\PY{n}{cm}\PY{o}{.}\PY{n}{deep}\PY{p}{,} \PY{o}{\PYZhy{}}\PY{l+m+mf}{0.20}\PY{p}{,} \PY{l+m+mf}{1.0}\PY{p}{)}\PY{p}{,}
    \PY{p}{(}\PY{l+s+s2}{\PYZdq{}}\PY{l+s+s2}{temperature}\PY{l+s+s2}{\PYZdq{}}\PY{p}{,} \PY{l+s+s2}{\PYZdq{}}\PY{l+s+s2}{Temperature (°C)}\PY{l+s+s2}{\PYZdq{}}\PY{p}{,} \PY{n}{cmocean}\PY{o}{.}\PY{n}{cm}\PY{o}{.}\PY{n}{thermal}\PY{p}{,} \PY{l+m+mf}{24.0}\PY{p}{,} \PY{l+m+mf}{34.0}\PY{p}{)}\PY{p}{,}
    \PY{p}{(}\PY{l+s+s2}{\PYZdq{}}\PY{l+s+s2}{salinity}\PY{l+s+s2}{\PYZdq{}}\PY{p}{,} \PY{l+s+s2}{\PYZdq{}}\PY{l+s+s2}{Salinity (PSU)}\PY{l+s+s2}{\PYZdq{}}\PY{p}{,} \PY{n}{cmocean}\PY{o}{.}\PY{n}{cm}\PY{o}{.}\PY{n}{haline}\PY{p}{,} \PY{l+m+mf}{0.0}\PY{p}{,} \PY{l+m+mf}{18.0}\PY{p}{)}\PY{p}{,}
\PY{p}{]}

\PY{c+c1}{\PYZsh{} Choose the timestep to preview.}
\PY{n}{preview\PYZus{}timestep} \PY{o}{=} \PY{l+m+mi}{0}
\PY{n}{t\PYZus{}str} \PY{o}{=} \PY{n}{time\PYZus{}str\PYZus{}all}\PY{o}{.}\PY{n}{isel}\PY{p}{(}\PY{n}{time}\PY{o}{=}\PY{n}{preview\PYZus{}timestep}\PY{p}{)}\PY{o}{.}\PY{n}{item}\PY{p}{(}\PY{p}{)}
\PY{n+nb}{print}\PY{p}{(}\PY{l+s+s2}{\PYZdq{}}\PY{l+s+s2}{Plot timestamp:}\PY{l+s+s2}{\PYZdq{}}\PY{p}{,} \PY{n}{t\PYZus{}str}\PY{p}{)}

\PY{c+c1}{\PYZsh{} Load UGRID dataset once and use renamed variables.}
\PY{n}{uds} \PY{o}{=} \PY{n}{xu}\PY{o}{.}\PY{n}{open\PYZus{}dataset}\PY{p}{(}\PY{n}{model\PYZus{}nc}\PY{p}{)}\PY{o}{.}\PY{n}{rename}\PY{p}{(}\PY{n}{var\PYZus{}name\PYZus{}map}\PY{p}{)}

\PY{c+c1}{\PYZsh{} Font\PYZhy{}size controls}
\PY{n}{title\PYZus{}fontsize} \PY{o}{=} \PY{l+m+mi}{9}
\PY{n}{axis\PYZus{}label\PYZus{}fontsize} \PY{o}{=} \PY{l+m+mi}{9}
\PY{n}{tick\PYZus{}fontsize} \PY{o}{=} \PY{l+m+mi}{8}
\PY{n}{cbar\PYZus{}label\PYZus{}fontsize} \PY{o}{=} \PY{l+m+mi}{9}
\PY{n}{cbar\PYZus{}tick\PYZus{}fontsize} \PY{o}{=} \PY{l+m+mi}{8}

\PY{c+c1}{\PYZsh{} Create Figure}
\PY{n}{fig}\PY{p}{,} \PY{n}{axes} \PY{o}{=} \PY{n}{plt}\PY{o}{.}\PY{n}{subplots}\PY{p}{(}\PY{l+m+mi}{2}\PY{p}{,} \PY{l+m+mi}{2}\PY{p}{,} \PY{n}{figsize}\PY{o}{=}\PY{p}{(}\PY{l+m+mi}{12}\PY{p}{,} \PY{l+m+mf}{5.5}\PY{p}{)}\PY{p}{,} \PY{n}{dpi}\PY{o}{=}\PY{l+m+mi}{200}\PY{p}{)}
\PY{n}{axes} \PY{o}{=} \PY{n}{axes}\PY{o}{.}\PY{n}{ravel}\PY{p}{(}\PY{p}{)}
\PY{k}{for} \PY{n}{ax}\PY{p}{,} \PY{p}{(}\PY{n}{var\PYZus{}name}\PY{p}{,} \PY{n}{common\PYZus{}name}\PY{p}{,} \PY{n}{cmap}\PY{p}{,} \PY{n}{vmin}\PY{p}{,} \PY{n}{vmax}\PY{p}{)} \PY{o+ow}{in} \PY{n+nb}{zip}\PY{p}{(}\PY{n}{axes}\PY{p}{,} \PY{n}{plot\PYZus{}specs}\PY{p}{)}\PY{p}{:}
    \PY{n}{quad} \PY{o}{=} \PY{n}{uds}\PY{p}{[}\PY{n}{var\PYZus{}name}\PY{p}{]}\PY{o}{.}\PY{n}{isel}\PY{p}{(}\PY{n}{time}\PY{o}{=}\PY{n}{preview\PYZus{}timestep}\PY{p}{)}\PY{o}{.}\PY{n}{ugrid}\PY{o}{.}\PY{n}{plot}\PY{p}{(}
        \PY{n}{ax}\PY{o}{=}\PY{n}{ax}\PY{p}{,}
        \PY{n}{cmap}\PY{o}{=}\PY{n}{cmap}\PY{p}{,}
        \PY{n}{vmin}\PY{o}{=}\PY{n}{vmin}\PY{p}{,}
        \PY{n}{vmax}\PY{o}{=}\PY{n}{vmax}\PY{p}{,}
        \PY{n}{add\PYZus{}colorbar}\PY{o}{=}\PY{k+kc}{True}\PY{p}{,}
        \PY{n}{cbar\PYZus{}kwargs}\PY{o}{=}\PY{p}{\PYZob{}}\PY{l+s+s2}{\PYZdq{}}\PY{l+s+s2}{label}\PY{l+s+s2}{\PYZdq{}}\PY{p}{:} \PY{n}{common\PYZus{}name}\PY{p}{\PYZcb{}}\PY{p}{,}
    \PY{p}{)}
    \PY{n}{cbar} \PY{o}{=} \PY{n}{quad}\PY{o}{.}\PY{n}{colorbar}
    \PY{n}{cbar}\PY{o}{.}\PY{n}{set\PYZus{}label}\PY{p}{(}\PY{n}{common\PYZus{}name}\PY{p}{,} \PY{n}{fontsize}\PY{o}{=}\PY{n}{cbar\PYZus{}label\PYZus{}fontsize}\PY{p}{)}
    \PY{n}{cbar}\PY{o}{.}\PY{n}{ax}\PY{o}{.}\PY{n}{tick\PYZus{}params}\PY{p}{(}\PY{n}{labelsize}\PY{o}{=}\PY{n}{cbar\PYZus{}tick\PYZus{}fontsize}\PY{p}{)}

    \PY{n}{ax}\PY{o}{.}\PY{n}{set\PYZus{}aspect}\PY{p}{(}\PY{l+s+s2}{\PYZdq{}}\PY{l+s+s2}{equal}\PY{l+s+s2}{\PYZdq{}}\PY{p}{)}
    \PY{n}{ax}\PY{o}{.}\PY{n}{set\PYZus{}title}\PY{p}{(}\PY{l+s+sa}{f}\PY{l+s+s2}{\PYZdq{}}\PY{l+s+s2}{Date: }\PY{l+s+si}{\PYZob{}}\PY{n}{t\PYZus{}str}\PY{l+s+si}{\PYZcb{}}\PY{l+s+s2}{\PYZdq{}}\PY{p}{,} \PY{n}{fontsize}\PY{o}{=}\PY{n}{title\PYZus{}fontsize}\PY{p}{)}
    \PY{n}{ax}\PY{o}{.}\PY{n}{set\PYZus{}xlabel}\PY{p}{(}\PY{l+s+s2}{\PYZdq{}}\PY{l+s+s2}{Longitude (°)}\PY{l+s+s2}{\PYZdq{}}\PY{p}{,} \PY{n}{fontsize}\PY{o}{=}\PY{n}{axis\PYZus{}label\PYZus{}fontsize}\PY{p}{)}
    \PY{n}{ax}\PY{o}{.}\PY{n}{set\PYZus{}ylabel}\PY{p}{(}\PY{l+s+s2}{\PYZdq{}}\PY{l+s+s2}{Latitude (°)}\PY{l+s+s2}{\PYZdq{}}\PY{p}{,} \PY{n}{fontsize}\PY{o}{=}\PY{n}{axis\PYZus{}label\PYZus{}fontsize}\PY{p}{)}
    \PY{n}{ax}\PY{o}{.}\PY{n}{set\PYZus{}aspect}\PY{p}{(}\PY{l+s+s2}{\PYZdq{}}\PY{l+s+s2}{equal}\PY{l+s+s2}{\PYZdq{}}\PY{p}{)}
    \PY{n}{ax}\PY{o}{.}\PY{n}{set\PYZus{}facecolor}\PY{p}{(}\PY{l+s+s2}{\PYZdq{}}\PY{l+s+s2}{whitesmoke}\PY{l+s+s2}{\PYZdq{}}\PY{p}{)}
    \PY{n}{ax}\PY{o}{.}\PY{n}{tick\PYZus{}params}\PY{p}{(}\PY{n}{labelsize}\PY{o}{=}\PY{n}{tick\PYZus{}fontsize}\PY{p}{)}

\PY{n}{plt}\PY{o}{.}\PY{n}{tight\PYZus{}layout}\PY{p}{(}\PY{p}{)}
\PY{n}{plt}\PY{o}{.}\PY{n}{subplots\PYZus{}adjust}\PY{p}{(}\PY{n}{hspace}\PY{o}{=}\PY{l+m+mf}{0.35}\PY{p}{)}
\PY{n}{plt}\PY{o}{.}\PY{n}{show}\PY{p}{(}\PY{p}{)}
\end{Verbatim}
\end{tcolorbox}

    \begin{Verbatim}[commandchars=\\\{\}]
Plot timestamp: 06-01-2013 00:00:00
    \end{Verbatim}

    \begin{center}
    \adjustimage{max size={0.9\linewidth}{0.9\paperheight}}{open_data_collaborator_code_notebook_files/open_data_collaborator_code_notebook_8_1.png}
    \end{center}
    { \hspace*{\fill} \\}
    
    \subsection{Plot salinity at one timestep}\label{plot-salinity-at-one-timestep}

This section makes a larger single-variable plot for salinity.


    \begin{tcolorbox}[breakable, size=fbox, boxrule=1pt, pad at break*=1mm,colback=cellbackground, colframe=cellborder]
\prompt{In}{incolor}{10}{\boxspacing}
\begin{Verbatim}[commandchars=\\\{\}]
\PY{n}{preview\PYZus{}timestep} \PY{o}{=} \PY{n+nb}{min}\PY{p}{(}\PY{l+m+mi}{12}\PY{p}{,} \PY{n}{n\PYZus{}timesteps} \PY{o}{\PYZhy{}} \PY{l+m+mi}{1}\PY{p}{)}
\PY{n}{t\PYZus{}str} \PY{o}{=} \PY{n}{time\PYZus{}str\PYZus{}all}\PY{o}{.}\PY{n}{isel}\PY{p}{(}\PY{n}{time}\PY{o}{=}\PY{n}{preview\PYZus{}timestep}\PY{p}{)}\PY{o}{.}\PY{n}{item}\PY{p}{(}\PY{p}{)}
\PY{n+nb}{print}\PY{p}{(}\PY{l+s+s2}{\PYZdq{}}\PY{l+s+s2}{Plot timestamp:}\PY{l+s+s2}{\PYZdq{}}\PY{p}{,} \PY{n}{t\PYZus{}str}\PY{p}{)}

\PY{n}{fig}\PY{p}{,} \PY{n}{ax} \PY{o}{=} \PY{n}{plt}\PY{o}{.}\PY{n}{subplots}\PY{p}{(}\PY{n}{figsize}\PY{o}{=}\PY{p}{(}\PY{l+m+mf}{6.2}\PY{p}{,} \PY{l+m+mf}{2.6}\PY{p}{)}\PY{p}{,} \PY{n}{dpi}\PY{o}{=}\PY{l+m+mi}{200}\PY{p}{)}
\PY{n}{quad} \PY{o}{=} \PY{n}{uds}\PY{p}{[}\PY{l+s+s2}{\PYZdq{}}\PY{l+s+s2}{salinity}\PY{l+s+s2}{\PYZdq{}}\PY{p}{]}\PY{o}{.}\PY{n}{isel}\PY{p}{(}\PY{n}{time}\PY{o}{=}\PY{n}{preview\PYZus{}timestep}\PY{p}{)}\PY{o}{.}\PY{n}{ugrid}\PY{o}{.}\PY{n}{plot}\PY{p}{(}
    \PY{n}{ax}\PY{o}{=}\PY{n}{ax}\PY{p}{,}
    \PY{n}{cmap}\PY{o}{=}\PY{n}{cmocean}\PY{o}{.}\PY{n}{cm}\PY{o}{.}\PY{n}{haline}\PY{p}{,}
    \PY{n}{vmin}\PY{o}{=}\PY{l+m+mf}{0.0}\PY{p}{,}
    \PY{n}{vmax}\PY{o}{=}\PY{l+m+mf}{18.0}\PY{p}{,}
    \PY{n}{add\PYZus{}colorbar}\PY{o}{=}\PY{k+kc}{True}\PY{p}{,}
    \PY{n}{cbar\PYZus{}kwargs}\PY{o}{=}\PY{p}{\PYZob{}}\PY{l+s+s2}{\PYZdq{}}\PY{l+s+s2}{label}\PY{l+s+s2}{\PYZdq{}}\PY{p}{:} \PY{l+s+s2}{\PYZdq{}}\PY{l+s+s2}{Salinity (PSU)}\PY{l+s+s2}{\PYZdq{}}\PY{p}{\PYZcb{}}\PY{p}{,}
\PY{p}{)}

\PY{n}{cbar} \PY{o}{=} \PY{n}{quad}\PY{o}{.}\PY{n}{colorbar}
\PY{n}{cbar}\PY{o}{.}\PY{n}{set\PYZus{}label}\PY{p}{(}\PY{l+s+s2}{\PYZdq{}}\PY{l+s+s2}{Salinity (PSU)}\PY{l+s+s2}{\PYZdq{}}\PY{p}{,} \PY{n}{fontsize}\PY{o}{=}\PY{l+m+mi}{9}\PY{p}{)}
\PY{n}{cbar}\PY{o}{.}\PY{n}{ax}\PY{o}{.}\PY{n}{tick\PYZus{}params}\PY{p}{(}\PY{n}{labelsize}\PY{o}{=}\PY{l+m+mi}{8}\PY{p}{)}

\PY{n}{ax}\PY{o}{.}\PY{n}{set\PYZus{}title}\PY{p}{(}\PY{l+s+sa}{f}\PY{l+s+s2}{\PYZdq{}}\PY{l+s+s2}{Date: }\PY{l+s+si}{\PYZob{}}\PY{n}{t\PYZus{}str}\PY{l+s+si}{\PYZcb{}}\PY{l+s+s2}{\PYZdq{}}\PY{p}{,} \PY{n}{fontsize}\PY{o}{=}\PY{l+m+mi}{9}\PY{p}{)}
\PY{n}{ax}\PY{o}{.}\PY{n}{set\PYZus{}xlabel}\PY{p}{(}\PY{l+s+s2}{\PYZdq{}}\PY{l+s+s2}{Longitude (°)}\PY{l+s+s2}{\PYZdq{}}\PY{p}{,} \PY{n}{fontsize}\PY{o}{=}\PY{l+m+mi}{9}\PY{p}{)}
\PY{n}{ax}\PY{o}{.}\PY{n}{set\PYZus{}ylabel}\PY{p}{(}\PY{l+s+s2}{\PYZdq{}}\PY{l+s+s2}{Latitude (°)}\PY{l+s+s2}{\PYZdq{}}\PY{p}{,} \PY{n}{fontsize}\PY{o}{=}\PY{l+m+mi}{9}\PY{p}{)}
\PY{n}{ax}\PY{o}{.}\PY{n}{set\PYZus{}aspect}\PY{p}{(}\PY{l+s+s2}{\PYZdq{}}\PY{l+s+s2}{equal}\PY{l+s+s2}{\PYZdq{}}\PY{p}{)}
\PY{n}{ax}\PY{o}{.}\PY{n}{set\PYZus{}facecolor}\PY{p}{(}\PY{l+s+s2}{\PYZdq{}}\PY{l+s+s2}{whitesmoke}\PY{l+s+s2}{\PYZdq{}}\PY{p}{)}
\PY{n}{ax}\PY{o}{.}\PY{n}{tick\PYZus{}params}\PY{p}{(}\PY{n}{labelsize}\PY{o}{=}\PY{l+m+mi}{8}\PY{p}{)}

\PY{n}{plt}\PY{o}{.}\PY{n}{tight\PYZus{}layout}\PY{p}{(}\PY{p}{)}
\PY{n}{plt}\PY{o}{.}\PY{n}{show}\PY{p}{(}\PY{p}{)}
\end{Verbatim}
\end{tcolorbox}

    \begin{Verbatim}[commandchars=\\\{\}]
Plot timestamp: 06-01-2013 12:00:00
    \end{Verbatim}

    \begin{center}
    \adjustimage{max size={0.9\linewidth}{0.9\paperheight}}{open_data_collaborator_code_notebook_files/open_data_collaborator_code_notebook_10_1.png}
    \end{center}
    { \hspace*{\fill} \\}
    
    \subsection{Create Video (optional)}\label{create-video-optional}

This section provides optional code for building an animation from the
NetCDF output.

The movie block is left commented so it does not run by accident. If you
want to use it:

\begin{itemize}
\tightlist
\item
  uncomment the lines in the code cell below
\item
  choose the variable to animate
\item
  set the timestep spacing and frames per second
\item
  run the block to save either an MP4 or, if ffmpeg is unavailable, a
  GIF
\end{itemize}


    \begin{tcolorbox}[breakable, size=fbox, boxrule=1pt, pad at break*=1mm,colback=cellbackground, colframe=cellborder]
\prompt{In}{incolor}{ }{\boxspacing}
\begin{Verbatim}[commandchars=\\\{\}]
\PY{c+c1}{\PYZsh{} Optional one\PYZhy{}time install for video export only.}
\PY{c+c1}{\PYZsh{} Keep this False for normal use.}
\PY{c+c1}{\PYZsh{} INSTALL\PYZus{}VIDEO\PYZus{}DEPS = False}
\PY{c+c1}{\PYZsh{} if INSTALL\PYZus{}VIDEO\PYZus{}DEPS:}
\PY{c+c1}{\PYZsh{}     \PYZsh{} Needed so Matplotlib can find ffmpeg for MP4 writing.}
\PY{c+c1}{\PYZsh{}     subprocess.check\PYZus{}call([sys.executable, \PYZdq{}\PYZhy{}m\PYZdq{}, \PYZdq{}pip\PYZdq{}, \PYZdq{}install\PYZdq{}, \PYZdq{}\PYZhy{}U\PYZdq{}, \PYZdq{}imageio\PYZhy{}ffmpeg\PYZdq{}])}

\PY{c+c1}{\PYZsh{} Choose which variable to animate: salinity, temperature, velocity, water\PYZus{}level}
\PY{c+c1}{\PYZsh{} movie\PYZus{}var = \PYZdq{}temperature\PYZdq{}  \PYZsh{} Change this to the variable you want to animate}

\PY{c+c1}{\PYZsh{} movie\PYZus{}specs = \PYZob{}}
\PY{c+c1}{\PYZsh{}     \PYZdq{}velocity\PYZdq{}: \PYZob{}}
\PY{c+c1}{\PYZsh{}         \PYZdq{}var\PYZus{}name\PYZdq{}: \PYZdq{}velocity\PYZus{}magnitude\PYZdq{},}
\PY{c+c1}{\PYZsh{}         \PYZdq{}label\PYZdq{}: \PYZdq{}Velocity Magnitude \PYZdl{}(m s\PYZca{}\PYZob{}\PYZhy{}1\PYZcb{})\PYZdl{}\PYZdq{},}
\PY{c+c1}{\PYZsh{}         \PYZdq{}cmap\PYZdq{}: cmocean.cm.amp,}
\PY{c+c1}{\PYZsh{}         \PYZdq{}vmin\PYZdq{}: 0.0,}
\PY{c+c1}{\PYZsh{}         \PYZdq{}vmax\PYZdq{}: 0.8,}
\PY{c+c1}{\PYZsh{}         \PYZdq{}file\PYZus{}tag\PYZdq{}: \PYZdq{}velocity\PYZdq{},}
\PY{c+c1}{\PYZsh{}     \PYZcb{},}
\PY{c+c1}{\PYZsh{}     \PYZdq{}water\PYZus{}level\PYZdq{}: \PYZob{}}
\PY{c+c1}{\PYZsh{}         \PYZdq{}var\PYZus{}name\PYZdq{}: \PYZdq{}water\PYZus{}level\PYZdq{},}
\PY{c+c1}{\PYZsh{}         \PYZdq{}label\PYZdq{}: \PYZdq{}Water Level (m)\PYZdq{},}
\PY{c+c1}{\PYZsh{}         \PYZdq{}cmap\PYZdq{}: cmocean.cm.deep,}
\PY{c+c1}{\PYZsh{}         \PYZdq{}vmin\PYZdq{}: \PYZhy{}0.20,}
\PY{c+c1}{\PYZsh{}         \PYZdq{}vmax\PYZdq{}: 1.0,}
\PY{c+c1}{\PYZsh{}         \PYZdq{}file\PYZus{}tag\PYZdq{}: \PYZdq{}water\PYZus{}level\PYZdq{},}
\PY{c+c1}{\PYZsh{}     \PYZcb{},}
\PY{c+c1}{\PYZsh{}     \PYZdq{}temperature\PYZdq{}: \PYZob{}}
\PY{c+c1}{\PYZsh{}         \PYZdq{}var\PYZus{}name\PYZdq{}: \PYZdq{}temperature\PYZdq{},}
\PY{c+c1}{\PYZsh{}         \PYZdq{}label\PYZdq{}: \PYZdq{}Temperature (°C)\PYZdq{},}
\PY{c+c1}{\PYZsh{}         \PYZdq{}cmap\PYZdq{}: cmocean.cm.thermal,}
\PY{c+c1}{\PYZsh{}         \PYZdq{}vmin\PYZdq{}: 24.0,}
\PY{c+c1}{\PYZsh{}         \PYZdq{}vmax\PYZdq{}: 34.0,}
\PY{c+c1}{\PYZsh{}         \PYZdq{}file\PYZus{}tag\PYZdq{}: \PYZdq{}temperature\PYZdq{},}
\PY{c+c1}{\PYZsh{}     \PYZcb{},}
\PY{c+c1}{\PYZsh{}     \PYZdq{}salinity\PYZdq{}: \PYZob{}}
\PY{c+c1}{\PYZsh{}         \PYZdq{}var\PYZus{}name\PYZdq{}: \PYZdq{}salinity\PYZdq{},}
\PY{c+c1}{\PYZsh{}         \PYZdq{}label\PYZdq{}: \PYZdq{}Salinity (PSU)\PYZdq{},}
\PY{c+c1}{\PYZsh{}         \PYZdq{}cmap\PYZdq{}: cmocean.cm.haline,}
\PY{c+c1}{\PYZsh{}         \PYZdq{}vmin\PYZdq{}: 0.0,}
\PY{c+c1}{\PYZsh{}         \PYZdq{}vmax\PYZdq{}: 18.0,}
\PY{c+c1}{\PYZsh{}         \PYZdq{}file\PYZus{}tag\PYZdq{}: \PYZdq{}salinity\PYZdq{},}
\PY{c+c1}{\PYZsh{}     \PYZcb{},}
\PY{c+c1}{\PYZsh{} \PYZcb{}}

\PY{c+c1}{\PYZsh{} if movie\PYZus{}var not in movie\PYZus{}specs:}
\PY{c+c1}{\PYZsh{}     raise ValueError(}
\PY{c+c1}{\PYZsh{}         f\PYZdq{}movie\PYZus{}var must be one of \PYZob{}list(movie\PYZus{}specs.keys())\PYZcb{}, got \PYZsq{}\PYZob{}movie\PYZus{}var\PYZcb{}\PYZsq{}.\PYZdq{}}
\PY{c+c1}{\PYZsh{}     )}

\PY{c+c1}{\PYZsh{} spec = movie\PYZus{}specs[movie\PYZus{}var]}
\PY{c+c1}{\PYZsh{} var\PYZus{}name = spec[\PYZdq{}var\PYZus{}name\PYZdq{}]}
\PY{c+c1}{\PYZsh{} label = spec[\PYZdq{}label\PYZdq{}]}
\PY{c+c1}{\PYZsh{} cmap = spec[\PYZdq{}cmap\PYZdq{}]}
\PY{c+c1}{\PYZsh{} vmin = spec[\PYZdq{}vmin\PYZdq{}]}
\PY{c+c1}{\PYZsh{} vmax = spec[\PYZdq{}vmax\PYZdq{}]}
\PY{c+c1}{\PYZsh{} file\PYZus{}tag = spec[\PYZdq{}file\PYZus{}tag\PYZdq{}]}

\PY{c+c1}{\PYZsh{} \PYZsh{} Try to point Matplotlib to bundled ffmpeg if imageio\PYZhy{}ffmpeg is installed.}
\PY{c+c1}{\PYZsh{} ffmpeg\PYZus{}ready = False}
\PY{c+c1}{\PYZsh{} try:}
\PY{c+c1}{\PYZsh{}     import imageio\PYZus{}ffmpeg}

\PY{c+c1}{\PYZsh{}     plt.rcParams[\PYZdq{}animation.ffmpeg\PYZus{}path\PYZdq{}] = imageio\PYZus{}ffmpeg.get\PYZus{}ffmpeg\PYZus{}exe()}
\PY{c+c1}{\PYZsh{}     ffmpeg\PYZus{}ready = True}
\PY{c+c1}{\PYZsh{} except Exception:}
\PY{c+c1}{\PYZsh{}     ffmpeg\PYZus{}ready = False}

\PY{c+c1}{\PYZsh{} \PYZsh{} Movie settings}
\PY{c+c1}{\PYZsh{} start\PYZus{}idx = 0}
\PY{c+c1}{\PYZsh{} end\PYZus{}idx = n\PYZus{}timesteps \PYZhy{} 1}
\PY{c+c1}{\PYZsh{} frame\PYZus{}step = 6          \PYZsh{} change this to control timestep spacing between frames}
\PY{c+c1}{\PYZsh{} fps = 6                 \PYZsh{} playback speed}
\PY{c+c1}{\PYZsh{} bitrate = 1800}

\PY{c+c1}{\PYZsh{} \PYZsh{} Build frame index list using the requested stride}
\PY{c+c1}{\PYZsh{} frame\PYZus{}indices = list(range(start\PYZus{}idx, end\PYZus{}idx + 1, frame\PYZus{}step))}
\PY{c+c1}{\PYZsh{} if not frame\PYZus{}indices:}
\PY{c+c1}{\PYZsh{}     raise ValueError(\PYZdq{}No frames selected. Check start\PYZus{}idx/end\PYZus{}idx/frame\PYZus{}step.\PYZdq{})}

\PY{c+c1}{\PYZsh{} \PYZsh{} Output location (next to the input NetCDF)}
\PY{c+c1}{\PYZsh{} video\PYZus{}dir = model\PYZus{}nc.parent}
\PY{c+c1}{\PYZsh{} video\PYZus{}dir.mkdir(parents=True, exist\PYZus{}ok=True)}
\PY{c+c1}{\PYZsh{} mp4\PYZus{}path = video\PYZus{}dir / f\PYZdq{}\PYZob{}file\PYZus{}tag\PYZcb{}\PYZus{}movie\PYZus{}step\PYZob{}frame\PYZus{}step\PYZcb{}\PYZus{}fps\PYZob{}fps\PYZcb{}.mp4\PYZdq{}}
\PY{c+c1}{\PYZsh{} gif\PYZus{}path = video\PYZus{}dir / f\PYZdq{}\PYZob{}file\PYZus{}tag\PYZcb{}\PYZus{}movie\PYZus{}step\PYZob{}frame\PYZus{}step\PYZcb{}\PYZus{}fps\PYZob{}fps\PYZcb{}.gif\PYZdq{}}

\PY{c+c1}{\PYZsh{} \PYZsh{} Create figure and first frame}
\PY{c+c1}{\PYZsh{} fig, ax = plt.subplots(figsize=(6.2, 2.6), dpi=200)}
\PY{c+c1}{\PYZsh{} first\PYZus{}idx = frame\PYZus{}indices[0]}
\PY{c+c1}{\PYZsh{} first\PYZus{}t = time\PYZus{}str\PYZus{}all.isel(time=first\PYZus{}idx).item()}
\PY{c+c1}{\PYZsh{} quad = uds[var\PYZus{}name].isel(time=first\PYZus{}idx).ugrid.plot(}
\PY{c+c1}{\PYZsh{}     ax=ax,}
\PY{c+c1}{\PYZsh{}     cmap=cmap,}
\PY{c+c1}{\PYZsh{}     vmin=vmin,}
\PY{c+c1}{\PYZsh{}     vmax=vmax,}
\PY{c+c1}{\PYZsh{}     add\PYZus{}colorbar=True,}
\PY{c+c1}{\PYZsh{}     cbar\PYZus{}kwargs=\PYZob{}\PYZdq{}label\PYZdq{}: label\PYZcb{},}
\PY{c+c1}{\PYZsh{} )}

\PY{c+c1}{\PYZsh{} cbar = quad.colorbar}
\PY{c+c1}{\PYZsh{} cbar.set\PYZus{}label(label, fontsize=9)}
\PY{c+c1}{\PYZsh{} cbar.ax.tick\PYZus{}params(labelsize=8)}

\PY{c+c1}{\PYZsh{} ax.set\PYZus{}xlabel(\PYZdq{}Longitude (°)\PYZdq{}, fontsize=9)}
\PY{c+c1}{\PYZsh{} ax.set\PYZus{}ylabel(\PYZdq{}Latitude (°)\PYZdq{}, fontsize=9)}
\PY{c+c1}{\PYZsh{} ax.set\PYZus{}aspect(\PYZdq{}equal\PYZdq{})}
\PY{c+c1}{\PYZsh{} ax.set\PYZus{}facecolor(\PYZdq{}whitesmoke\PYZdq{})}
\PY{c+c1}{\PYZsh{} ax.tick\PYZus{}params(labelsize=8)}
\PY{c+c1}{\PYZsh{} ax.set\PYZus{}title(f\PYZdq{}Date: \PYZob{}first\PYZus{}t\PYZcb{}\PYZdq{}, fontsize=9)}


\PY{c+c1}{\PYZsh{} def update(frame\PYZus{}num):}
\PY{c+c1}{\PYZsh{}     idx = frame\PYZus{}indices[frame\PYZus{}num]}
\PY{c+c1}{\PYZsh{}     t\PYZus{}str = time\PYZus{}str\PYZus{}all.isel(time=idx).item()}

\PY{c+c1}{\PYZsh{}     \PYZsh{} Remove previous mesh but keep colorbar/axes styling intact.}
\PY{c+c1}{\PYZsh{}     if ax.collections:}
\PY{c+c1}{\PYZsh{}         ax.collections[0].remove()}

\PY{c+c1}{\PYZsh{}     uds[var\PYZus{}name].isel(time=idx).ugrid.plot(}
\PY{c+c1}{\PYZsh{}         ax=ax,}
\PY{c+c1}{\PYZsh{}         cmap=cmap,}
\PY{c+c1}{\PYZsh{}         vmin=vmin,}
\PY{c+c1}{\PYZsh{}         vmax=vmax,}
\PY{c+c1}{\PYZsh{}         add\PYZus{}colorbar=False,}
\PY{c+c1}{\PYZsh{}     )}
\PY{c+c1}{\PYZsh{}     \PYZsh{} Re\PYZhy{}apply labels and styling because ugrid.plot can reset axis labels to x/y.}
\PY{c+c1}{\PYZsh{}     ax.set\PYZus{}xlabel(\PYZdq{}Longitude (°)\PYZdq{}, fontsize=9)}
\PY{c+c1}{\PYZsh{}     ax.set\PYZus{}ylabel(\PYZdq{}Latitude (°)\PYZdq{}, fontsize=9)}
\PY{c+c1}{\PYZsh{}     ax.set\PYZus{}aspect(\PYZdq{}equal\PYZdq{})}
\PY{c+c1}{\PYZsh{}     ax.set\PYZus{}facecolor(\PYZdq{}whitesmoke\PYZdq{})}
\PY{c+c1}{\PYZsh{}     ax.tick\PYZus{}params(labelsize=8)}
\PY{c+c1}{\PYZsh{}     ax.set\PYZus{}title(f\PYZdq{}Date: \PYZob{}t\PYZus{}str\PYZcb{}\PYZdq{}, fontsize=9)}
\PY{c+c1}{\PYZsh{}     return ax.collections}


\PY{c+c1}{\PYZsh{} anim = FuncAnimation(}
\PY{c+c1}{\PYZsh{}     fig,}
\PY{c+c1}{\PYZsh{}     update,}
\PY{c+c1}{\PYZsh{}     frames=len(frame\PYZus{}indices),}
\PY{c+c1}{\PYZsh{}     interval=1000 / fps,}
\PY{c+c1}{\PYZsh{}     blit=False,}
\PY{c+c1}{\PYZsh{} )}

\PY{c+c1}{\PYZsh{} if ffmpeg\PYZus{}ready:}
\PY{c+c1}{\PYZsh{}     writer = FFMpegWriter(fps=fps, bitrate=bitrate)}
\PY{c+c1}{\PYZsh{}     anim.save(mp4\PYZus{}path, writer=writer)}
\PY{c+c1}{\PYZsh{}     saved\PYZus{}path = mp4\PYZus{}path}
\PY{c+c1}{\PYZsh{} else:}
\PY{c+c1}{\PYZsh{}     \PYZsh{} Simple fallback when ffmpeg is unavailable.}
\PY{c+c1}{\PYZsh{}     writer = PillowWriter(fps=fps)}
\PY{c+c1}{\PYZsh{}     anim.save(gif\PYZus{}path, writer=writer)}
\PY{c+c1}{\PYZsh{}     saved\PYZus{}path = gif\PYZus{}path}
\PY{c+c1}{\PYZsh{}     print(\PYZdq{}ffmpeg not found. Saved GIF instead of MP4.\PYZdq{})}
\PY{c+c1}{\PYZsh{}     print(\PYZdq{}Set INSTALL\PYZus{}VIDEO\PYZus{}DEPS = True in this notebook block, run once, then rerun for MP4.\PYZdq{})}

\PY{c+c1}{\PYZsh{} plt.close(fig)}

\PY{c+c1}{\PYZsh{} print(f\PYZdq{}Saved \PYZob{}movie\PYZus{}var\PYZcb{} movie: \PYZob{}saved\PYZus{}path\PYZcb{}\PYZdq{})}
\PY{c+c1}{\PYZsh{} print(f\PYZdq{}Frames: \PYZob{}len(frame\PYZus{}indices)\PYZcb{} | Step: \PYZob{}frame\PYZus{}step\PYZcb{} timestep(s) | FPS: \PYZob{}fps\PYZcb{}\PYZdq{})}
\end{Verbatim}
\end{tcolorbox}


    % Add a bibliography block to the postdoc
    
    
    
\end{document}
